\section{System Actors}
The roles and actors defined in the Neo\ac{LAAC} ecosystem are described in the table below:

\vspace{0.3cm}
\begin{tabular}{|l|p{11.5cm}|}
\hline
\textbf{Actor} & \textbf{Description} \\ \hline
\textbf{Student / Freshmen} & Principal end user. Consults academic information, maps, public transport routes, interacts in the social feed, and uses the chat and panic/emergency button channels. \\ \hline
\textbf{Professor / Faculty} & User responsible for sharing announcements, managing academic support hours, and interacting academically with students. \\ \hline
\textbf{Staff \ac{LAAC} - Frontdesk} & Members of the real-time physical/human support team. Responds to incidents and high-priority tickets escalated by the chatbot. \\ \hline
\textbf{Staff \ac{LAAC} - Devs / Testers} & System administrators and technical staff. Perform development changes, testing, and monitor audit logs. \\ \hline
\textbf{Staff \ac{LAAC} - Marketing} & Focused on promotion of communication channels, managing partnerships, and local discounts. \\ \hline
\textbf{Academic Organizations} & Student branches, unions, or student groups. Manage their own events and announcements in their "Nucleus" space. \\ \hline
\textbf{Local Partners} & External commercial entities that offer discounts and benefits to students. \\ \hline
\end{tabular}

\vspace{0.6cm}

\section{Microservices and Components Mapping}
The system is built on an architecture of independent containerized microservices:
\begin{itemize}
    \item \textbf{\texttt{frontend}}: Responsive web interface served by Nginx.
    \item \textbf{\texttt{orq}}: \ac{API} Orchestrator (Single entry Gateway for external communication).
    \item \textbf{\texttt{auth-service} / \texttt{profile-service}}: Handles user identity, registration, login, profile data, and \ac{RBAC} permissions.
    \item \textbf{\texttt{map-service}}: Integrates the \ac{UBI} campus map, geolocation, and pedestrian routing.
    \item \textbf{\texttt{faq-service} / \texttt{chatbot-service}}: \ac{FAQ} database and intelligent automated triage assistant.
    \item \textbf{\texttt{chat-service} / \texttt{ticket-service} / \texttt{notification-service}}: Real-time messaging, support agent workspace, and notifications dispatch.
    \item \textbf{\texttt{emergency-service} / \texttt{emergencyCall-service}}: Handles the panic button and phone callback requests.
    \item \textbf{\texttt{academic-service} / \texttt{calendar-service}}: Courses, subjects, schedules, class registration, and semester calendar sync.
    \item \textbf{\texttt{post-service} / \texttt{feed-service} / \texttt{events-service} / \texttt{news-service}}: Core community features (social feed, posts, likes, comments, news, and campus events).
    \item \textbf{\texttt{backup-service} / \texttt{logging-service}}: Security audits, logs central aggregation, and disaster recovery dumps.
\end{itemize}

\newpage

\section{Functional Requirements (FR)}
The functional requirements describe the actions and behaviors of the Neo\ac{LAAC} platform.

\subsection{Identity and Profile Management}
\begin{tabular}{|l|l|p{7.5cm}|l|}
\hline
\textbf{ID} & \textbf{Name} & \textbf{Description} & \textbf{Component} \\ \hline
\textbf{FR01} & User Registration & Allow users to register accounts using institutional email, password, and full name. & \texttt{auth-service} \\ \hline
\textbf{FR02} & Secure Auth & Authenticate users and generate secure \ac{JWT} tokens for session control. & \texttt{auth-service} \\ \hline
\textbf{FR03} & Dynamic Profiles & Support different profile views and fields based on user role (Student, Faculty, Staff). & \texttt{profile-service} \\ \hline
\textbf{FR04} & Profile Updates & Let users update personal info and upload avatar pictures. & \texttt{profile-service} \\ \hline
\textbf{FR05} & Roles Management & Create and modify roles and access levels administratively. & \texttt{auth-service} \\ \hline
\textbf{FR06} & Social Following & Allow users to follow and unfollow other academic profiles. & \texttt{profile-service} \\ \hline
\end{tabular}

\subsection{Orientation, Infrastructure, and Maps}
\begin{tabular}{|l|l|p{7.5cm}|l|}
\hline
\textbf{ID} & \textbf{Name} & \textbf{Description} & \textbf{Component} \\ \hline
\textbf{FR07} & \ac{UBI} Campus Map & Display an interactive map with campus buildings, floors, and classroom codes. & \texttt{map-service} \\ \hline
\textbf{FR08} & Pedestrian Routing & Calculate a path from current \ac{GPS} coordinates to a target classroom, supporting time offsets. & \texttt{map-service} \\ \hline
\textbf{FR09} & Transit Schedules & Detail municipal public transport schedules, routes, and pricing. & \texttt{map-service} \\ \hline
\end{tabular}

\subsection{Hybrid Support, Emergencies, and Tickets}
\begin{tabular}{|l|l|p{7.5cm}|l|}
\hline
\textbf{ID} & \textbf{Name} & \textbf{Description} & \textbf{Component} \\ \hline
\textbf{FR10} & \ac{FAQ} Section & Provide a list of categorized frequently asked questions. & \texttt{faq-service} \\ \hline
\textbf{FR11} & Chatbot Support & Enable users to ask questions in natural language to a virtual assistant ("Mentor"). & \texttt{chatbot-service} \\ \hline
\textbf{FR12} & Hybrid Triage & Monitor chat conversations and automatically transfer to human Frontdesk staff upon detecting danger keywords or bot failures. & \texttt{orq} / \texttt{chat-service} \\ \hline
\textbf{FR13} & Support Alerts & Dispatch push notifications to support agents when a chat session is escalated. & \texttt{notification-service} \\ \hline
\textbf{FR14} & Chat History & Persist and retrieve chat logs chronologically for students and agents. & \texttt{chat-service} \\ \hline
\textbf{FR15} & Callback Request & Allow users to request an immediate phone call callback in case of issues. & \texttt{emergencyCall-service} \\ \hline
\textbf{FR16} & Panic Button & Emit critical emergency alerts indicating geolocation to the security team. & \texttt{emergency-service} \\ \hline
\textbf{FR17} & Ticket Queue & Create and track the lifecycle of help tickets assigned to staff teams. & \texttt{ticket-service} \\ \hline
\end{tabular}

\subsection{Academic, Social Life, and Utilities}
\begin{tabular}{|l|l|p{7.5cm}|l|}
\hline
\textbf{ID} & \textbf{Name} & \textbf{Description} & \textbf{Component} \\ \hline
\textbf{FR18} & Courses Catalog & List academic courses and respective syllabi and professors. & \texttt{academic-service} \\ \hline
\textbf{FR19} & Class Schedule & Show student schedules dynamically based on course enrollments. & \texttt{calendar-service} \\ \hline
\textbf{FR20} & Academic Calendar & Display general semester dates, holidays, and exam periods. & \texttt{calendar-service} \\ \hline
\textbf{FR21} & Schedule Reports & Allow students and professors to report scheduling conflicts. & \texttt{calendar-service} \\ \hline
\textbf{FR23} & Post Creation & Let users and organizations publish posts with text, images, and videos. & \texttt{post-service} \\ \hline
\textbf{FR25} & Integrated Feed & Aggregate campus announcements, general news, and community posts. & \texttt{feed-service} \\ \hline
\textbf{FR29} & Central Audits & Aggregate and expose service logs to authorized administrators. & \texttt{logging-service} \\ \hline
\textbf{FR30} & Scheduled Backups & Export database dumps to the backup volume at regular intervals. & \texttt{backup-service} \\ \hline
\end{tabular}

\newpage

\section{Non-Functional Requirements (NFR)}

\subsection{Security and Privacy}
\begin{itemize}
    \item \textbf{NFR01 (Encryption):} Secure password hashing using the \texttt{bcrypt} one-way algorithm.
    \item \textbf{NFR02 (Isolation):} Public port exposure restricted strictly to the Nginx \texttt{frontend} (8080/80). Microservices communicate entirely on private networks.
    \item \textbf{NFR03 (Access Control):} Administrative routes require \ac{JWT} authentication and Role-Based Access Control (\ac{RBAC}) validation.
    \item \textbf{NFR04 (File Security):} Access restrictions implemented for uploaded files ("public", "private", "academic-only").
    \item \textbf{NFR04a (Right to be Forgotten - \ac{GDPR}):} Provide a mechanism to permanently delete user accounts, deleting profile data and anonymizing support tickets (replacing info with "Deleted User").
    \item \textbf{NFR04b (Explicit Consent - \ac{GDPR}):} Explicit consent must be collected during registration for handling profile data, logs, and geolocation.
    \item \textbf{NFR04c (Retention Policies - \ac{GDPR}):}
    \begin{itemize}
        \item Audit logs: Retained for a maximum period of 90 days.
        \item Chat history and tickets: Archived/anonymized after 1 academic year.
        \item Geolocation: Processed strictly in-memory during active route calculation; persistency on disk is prohibited.
    \end{itemize}
\end{itemize}

\subsection{Performance and Scalability}
\begin{itemize}
    \item \textbf{NFR05 (Response Time):} Orchestrator response time must be under 200ms for read requests.
    \item \textbf{NFR06 (Chat Concurrency):} Real-time messaging delivery latency must remain below 1 second.
    \item \textbf{NFR07 (Scalability):} Microservices packaged as Docker images for rapid horizontal autoscaling in Kubernetes.
    \item \textbf{NFR08 (Databases):} Logical partition between \texttt{laac-mariadb} and \texttt{laac-calendar-mariadb} databases to prevent \ac{IO} bottlenecks during peak class search hours.
\end{itemize}

\subsection{Reliability, Resilience, and Availability}
\begin{itemize}
    \item \textbf{NFR09 (Fallback Mechanism):} If \texttt{chatbot-service} times out or throws an error, the orchestrator must instantly route the user to a live support agent and display a clear notification.
    \item \textbf{NFR10 (Health Checks):} Expose a \texttt{/health} endpoint in all services to support automated health monitoring and container restarts by Kubernetes.
    \item \textbf{NFR11 (Backups):} Automated backups daily with \ac{RPO} (Recovery Point Objective) under 24h and \ac{RTO} (Recovery Time Objective) under 2h.
\end{itemize}

\subsection{Usability, Accessibility, and Compatibility}
\begin{itemize}
    \item \textbf{NFR12 (Responsiveness):} 100\% mobile-first responsive layout tested across Chrome, Safari, Edge, and Firefox.
    \item \textbf{NFR13 (Language):} Native interface support in Portuguese (PT-PT).
    \item \textbf{NFR13a (Digital Accessibility - \ac{WCAG} 2.1 AA):} Frontend must strictly conform to \ac{WCAG} 2.1 Level AA specifications.
    \item \textbf{NFR13b (Screen Readers):} Semantically structured \ac{HTML} elements, plus \texttt{aria-label} and \texttt{alt} attributes, supporting NVDA/JAWS readers.
    \item \textbf{NFR13c (Keyboard-only):} Full \ac{UI} navigation using only the keyboard (\texttt{Tab} key navigation with distinct focus outlines).
    \item \textbf{NFR13d (Contrast Ratio):} Contrast ratio of at least 4.5:1 for normal text and 3:1 for large text.
\end{itemize}

\newpage

\section{Triage and Escalation Flow (Chatbot $\rightarrow$ Human)}
The chat triage pipeline model represents the automated decision tree when processing messages sent by students, determining when to answer via the virtual assistant or escalate to live support staff.

\begin{figure}[htbp]
    \centering
    \includegraphics[width=0.9\textwidth]{images/triage_flow.png}
    \caption{Triage and escalation decision pipeline diagram showing automated keywords checks, fallback routing, and live agent notification flow.}
    \label{fig:triage_flow}
\end{figure}

\newpage

\section{Requirements and Product Engineering}

\subsection{System Personas}
\begin{itemize}
    \item \textbf{Tiago Silva (Freshman):} 18 years old, freshman studying Computer Science, moved from Portimão. Needs to locate classroom block VI and find local bus times. Feels anxious and overwhelmed.
    \item \textbf{Dra. Maria Antunes (Faculty):} 45 years old, CS course coordinator. Needs to notify students of room shifts and upload office hours. Feels frustrated with students ignoring academic emails.
    \item \textbf{Rita Gomes (Frontdesk):} 21 years old, Psychology student and welcomer volunteer. Needs to filter out redundant questions to focus on urgent support requests.
\end{itemize}

\subsection{User Journey}
Tiago Silva's user journey on Day 1:
\begin{enumerate}
    \item Wakes up worried about bus lines ➔ Opens Neo\ac{LAAC} and finds local bus transit timetables (\textbf{\ac{MS}\_Map}).
    \item Arrives at campus, cannot find POO room ➔ Calculates pedestrian route to Block VI room (\textbf{\ac{MS}\_Map}).
    \item Has questions about academic integration guidelines ➔ Asks the Mentor Bot (\textbf{\ac{MS}\_Bot}).
    \item Bot fails to answer a complex, specific question ➔ Automatically transfers to Rita Gomes (\textbf{\ac{MS}\_Chat}).
\end{enumerate}

\subsection{Epics and User Stories}
\begin{itemize}
    \item \textbf{EPIC 01 - Spatial and Academic Orientation}:
    \begin{itemize}
        \item \textbf{\ac{US}01:} As a student, I want to view my consolidated class schedule to plan my day.
        \item \textbf{\ac{US}02:} As a student, I want step-by-step pedestrian routes to classrooms to avoid getting lost.
    \end{itemize}
    \item \textbf{EPIC 02 - Hybrid Support and Triage}:
    \begin{itemize}
        \item \textbf{\ac{US}03:} As a student, I want to chat with a bot to resolve common integration doubts instantly.
        \item \textbf{\ac{US}04:} As a student in danger, I want to be connected to a human immediately for support.
    \end{itemize}
\end{itemize}

\subsection{Acceptance Criteria}
Example for \textbf{\ac{US}04 (Hybrid Triage)}:
\begin{itemize}
    \item \textbf{Scenario 1 (Critical Keywords):} If the message contains "emergency", "panic", or "accident", the orchestrator changes the status to \texttt{escalated\_to\_human}, fires a push notification to Frontdesk, and informs the student.
    \item \textbf{Scenario 2 (Bot Failure):} If the bot returns the default fallback error phrase, the system logs the ticket and escalates to human live support.
\end{itemize}

\subsection{\ac{UML} Use Cases}
\begin{enumerate}
    \item \textbf{\ac{UC}01 - Consult Schedule:} Student requests schedule; \texttt{calendar-service} returns academic calendar.
    \item \textbf{\ac{UC}02 - Get Route:} Student requests route; \texttt{map-service} returns \ac{GeoJSON} path.
    \item \textbf{\ac{UC}03 - Get Support:} Student sends chat; \texttt{chatbot-service} processes; if needed, \texttt{orq} routes to live staff.
\end{enumerate}

\subsection{Product Backlog (MoSCoW)}
\begin{itemize}
    \item \textbf{Must Have:} Login/\ac{JWT}, Chatbot + Hybrid Triage, Panic Button.
    \item \textbf{Should Have:} Map + Pedestrian Routes, Class Schedules, Tickets Dashboard.
    \item \textbf{Could Have:} Social Feed, Posts/Likes, Local Partner discounts list.
    \item \textbf{Won't Have:} Video calls (text chat and phone callback are sufficient).
\end{itemize}

\newpage

\section{Design and Interface (\ac{UI}/\ac{UX})}

\subsection{Sitemap}
\begin{verbatim}
[Frontend Home]
 ├── /login (Auth and Signup)
 ├── /dashboard (Announcements and Welcome Board)
 ├── /map (Campus Map and Routing)
 ├── /chat (Mentor Bot / Panic Button)
 ├── /schedules (Class Schedules)
 └── /community (Social Feed and Partner Discounts)
\end{verbatim}

\subsection{Design System}
\begin{itemize}
    \item \textbf{Colors:} \ac{UBI} Blue (\texttt{\#005691}), Dark Cobalt (\texttt{\#002D62}), Success Green (\texttt{\#2E7D32}), Alarm Red (\texttt{\#D32F2F}).
    \item \textbf{Typography:} \texttt{Inter} (\ac{UI} body) and \texttt{Outfit} (headings).
    \item \textbf{Components:} \texttt{ActionButton} (rounded buttons with 8px radius and scale animations on hover); \texttt{ChatBubble} (student messages on the right in blue, system/bot responses on the left in gray).
\end{itemize}

\newpage

\section{Detailed Modeling and Architecture}

\subsection{Sequence Diagram (Hybrid Triage)}
The dynamic interaction sequence between system components during an emergency event is detailed in Figure~\ref{fig:sequence_diagram}. It outlines the lifecycle of a critical message from the initial client event to backend persistence, threat detection, and push notifications.

\begin{figure}[htbp]
    \centering
    \includegraphics[width=0.9\textwidth]{images/sequence_diagram.png}
    \caption{UML Sequence Diagram detailing message routing, synchronous database logs persistence, and threat detection escalation events inside the Orchestrator.}
    \label{fig:sequence_diagram}
\end{figure}

\subsection{OpenAPI Specification (Route Calculation)}
\begin{verbatim}
openapi: 3.0.3
info:
  title: NeoLAAC Map Service API
  version: 1.0.0
paths:
  /route:
    get:
      summary: Calculates pedestrian route to a campus location
      parameters:
        - name: start_lat
          in: query
          required: true
          schema:
            type: number
            format: float
        - name: start_lng
          in: query
          required: true
          schema:
            type: number
            format: float
        - name: dest_id
          in: query
          required: true
          schema:
            type: string
      responses:
        '200':
          description: Route calculated successfully
          content:
            application/json:
              schema:
                type: object
                properties:
                  distance_meters:
                    type: number
                  polyline_geojson:
                    type: string
\end{verbatim}

\subsection{Entity-Relationship Diagram (\ac{ERD}) and Cross-Service Management}
The data schema layout across independent service boundaries is shown in Figure~\ref{fig:erd_diagram}. It highlights logical key associations that span across multiple decoupled physical databases.

\begin{figure}[htbp]
    \centering
    \includegraphics[width=0.9\textwidth]{images/erd_diagram.png}
    \caption{Logical Entity-Relationship Diagram detailing cross-service data relations, primary keys, and logical foreign key constraints handled by the API Orchestrator.}
    \label{fig:erd_diagram}
\end{figure}

\begin{itemize}
    \item \textbf{Database Segregation:} The schedules microservice uses the independent database \texttt{laac-calendar-mariadb}. The relationship to the user is logical (\texttt{student\_id} references \texttt{users.id} in the main database \texttt{laac-mariadb}).
    \item \textbf{Integrity Management:} No hard foreign keys exist at the database network boundary. The \ac{API} Orchestrator (\texttt{orq}) validates relations synchronously using \ac{HTTP} endpoints.
\end{itemize}

\subsection{Ticket Lifecycle (State Machine)}
The transitions between support ticket statuses are formalized in Figure~\ref{fig:ticket_state_machine}, specifying validation rules for agent assignments and student contest triggers.

\begin{figure}[htbp]
    \centering
    \includegraphics[width=0.9\textwidth]{images/ticket_state_machine.png}
    \caption{State Machine Diagram illustrating ticket status transitions from opening, execution assignment, cancellation, resolution, and subsequent reopen windows.}
    \label{fig:ticket_state_machine}
\end{figure}

\newpage

\section{Development Patterns and Code}

\subsection{Design Patterns}
\begin{enumerate}
    \item \textbf{\ac{API} Gateway / Orchestrator}: Access orchestration and composition consolidated into \texttt{laac-orq}.
    \item \textbf{Database-per-Service}: Domain boundaries isolated using distinct MariaDB container nodes.
    \item \textbf{Command / Dispatcher Pattern}: Used by the Gateway to delegate requests based on message classification.
\end{enumerate}

\subsection{Commit Pattern Guide (Conventional Commits)}
\begin{itemize}
    \item \texttt{feat(scope):} New feature.
    \item \texttt{fix(scope):} Bug resolution.
    \item \texttt{chore(scope):} Dev environment updates and Dockerfiles.
\end{itemize}

\subsection{Git Branching Strategy}
The development lifecycle conforms to \textbf{GitFlow}:
\begin{itemize}
    \item \textbf{\texttt{main}}: Stable production branch. Protected. Only merges from \texttt{develop} (via \ac{PR}) and \texttt{hotfix/*} branches allowed.
    \item \textbf{\texttt{develop}}: Active integration branch.
    \item \textbf{\texttt{feature/*}}: Branches created off \texttt{develop} to implement product stories.
    \item \textbf{\texttt{hotfix/*}}: Branches branched off \texttt{main} for hot-patching critical bugs.
\end{itemize}

\newpage

\section{Quality and Testing (\ac{QA})}

\subsection{Test Coverage per Microservice}
\begin{tabular}{|l|c|l|p{7.5cm}|}
\hline
\textbf{Component} & \textbf{Tests} & \textbf{File} & \textbf{Validated Business Logic} \\ \hline
\texttt{orq} & 25 & \texttt{orq/tests/test\_orq.py} & Gateway routing, chat triage, and escalation. \\ \hline
\texttt{auth-service} & 12 & \texttt{auth/tests/test\_auth.py} & Login, signup, \ac{RBAC} levels, and \ac{JWT} validation. \\ \hline
\texttt{profile-service} & 11 & \texttt{profile/tests/test\_profile.py} & User profiles, follow toggle, and avatar uploads. \\ \hline
\texttt{emergency-service} & 5 & \texttt{emergency/tests/test\_emergency.py} & Alerter, history logs, and notifications trigger. \\ \hline
\texttt{ticket-service} & 3 & \texttt{ticket/tests/test\_main.py} & Ticket structure data validations. \\ \hline
\end{tabular}

\subsection{Hybrid Test Pipeline}
The \ac{CI}/\ac{CD} pipeline within \texttt{pipeline.py} operates with dynamic test discovery:
\begin{enumerate}
    \item Tries to execute the test suite using \texttt{pytest}.
    \item If \texttt{pytest} is missing from the container image, falls back to the native Python test runner:
    \begin{verbatim}
    python -m unittest discover -s tests
    \end{verbatim}
\end{enumerate}

\subsection{Test Script Example (Unittest)}
\begin{verbatim}
import unittest
from unittest.mock import patch, MagicMock
from fastapi.testclient import TestClient
from main import app, get_db

class TestAuthService(unittest.TestCase):
    def setUp(self):
        self.client = TestClient(app)
        self.mock_db = MagicMock()
        app.dependency_overrides[get_db] = lambda: self.mock_db

    def test_register_new_user(self):
        self.mock_db.execute.return_value.fetchone.return_value = None
        response = self.client.post("/register", json={
            "email": "new@ubi.pt", "password": "pass", "role": "student"
        })
        self.assertEqual(response.status_code, 200)
\end{verbatim}

\newpage

\section{Agile Project Management}
\begin{itemize}
    \item \textbf{Sprints}: 2-week iteration cycles. Includes planning, daily syncs, and retrospectives.
    \item \textbf{Definition of Done (\ac{DoD})}: Linter checks passed, code coverage $\ge$ 80\%, Docker build successful, and \ac{PR} peer-reviewed by at least one engineer.
\end{itemize}

\newpage

\section{DevOps, Deployment, and DevSecOps}

\subsection{Physical Deployment (Kubernetes Topology)}
The infrastructure setup in production environment is represented in Figure~\ref{fig:kubernetes_deployment}, showing the load balancing, network separation, pod distribution, and persistent volumes.

\begin{figure}[htbp]
    \centering
    \includegraphics[width=0.9\textwidth]{images/kubernetes_deployment.png}
    \caption{Infrastructure deployment architecture detailing Kubernetes Ingress mapping, decoupled service pods, cluster network boundaries, and database storage mappings.}
    \label{fig:kubernetes_deployment}
\end{figure}

\subsection{Orchestrator Dockerfile}
\begin{verbatim}
FROM python:3.11-slim
WORKDIR /app
RUN apt-get update && apt-get install -y --no-install-recommends \
    gcc libmariadb-dev && rm -rf /var/lib/apt/lists/*
COPY requirements.txt .
RUN pip install --no-cache-dir -r requirements.txt
COPY . .
EXPOSE 8000
CMD ["uvicorn", "main:app", "--host", "0.0.0.0", "--port", "8000"]
\end{verbatim}

\subsection{\ac{CI}/\ac{CD} Workflow (GitHub Actions)}
\begin{verbatim}
name: NeoLAAC Continuous Integration
on:
  push:
    branches: [ main, develop ]
jobs:
  validate:
    runs-on: ubuntu-latest
    steps:
    - uses: actions/checkout@v3
    - name: Set up Python
      uses: actions/setup-python@v4
      with:
        python-version: '3.11'
    - name: Run Pipeline Validation Script
      run: python pipeline.py
\end{verbatim}

\subsection{DevSecOps}
\begin{itemize}
    \item \textbf{Static Analysis (\ac{SAST})}: \texttt{Bandit} scanner integrated to find vulnerabilities in Python files.
    \item \textbf{Secrets Isolation}: All credentials injected at runtime via environment variables using a secure \texttt{.env} file.
\end{itemize}
